\section{Objetivos}

Este trabalho, de forma geral, visa a construção de uma extensão protótipo para navegadores Google Chrome, denominada ESPEON. A ferramenta é crucial para a coleta de dados comportamentais dos usuários submetidos a aulas expositivas em ambiente remoto de ensino, permitindo a identificação de dispersão de foco e falta de engajamento. Como objetivos específicos, a pesquisa busca avaliar o desempenho da metodologia expostivia de ensino, aplicada ao ensino remoto; tendo a análise pautada em indicadores de dispersão de foco e falta de engajamento, emitidos em relatórios pela extensão ESPEON. 